\chapter{Conclusion}
\label{chap:11-conclusion}

\section{La réponse}
\label{sec:reponse}

\begin{keybox}
Joelle Pineau a été écoutée sur ce qui accompagne le résultat : le lien vers le code, le calcul
chiffré, l'effet d'un réglage et, de façon plus fragile, la barre d'erreur. La barre s'est répandue sans que l'incertitude
calculée sur des graines aléatoires la suive de façon détectable. Le réglage des méthodes de
référence ne montre aucun signe de progrès.
\end{keybox}

Cette réponse concerne le RL classique. Des agents IA ont lu cinquante articles de chaque côté.
Nous avons ensuite comparé à notre propre lecture de 2018, toujours avec les mêmes grilles
(\cref{chap:03-sondage-methode,chap:04-sondage-resultats}). Nous avons de plus refait des
expériences avec les algorithmes de 2026
(\cref{chap:05-experiences-methode,chap:06-experiences-resultats,chap:07-modeles-de-langue}). Les
chiffres qui suivent sont issus de ces chapitres. Nous nous bornons ici à résumer ce que chacun
établit et renvoyons au chapitre concerné.

\section{Ce que les mesures établissent}
\label{sec:nombres}

\paragraph{Quatre pratiques ont changé} Le lien vers le code, le calcul chiffré, la sensibilité à un
réglage et l'incertitude rapportée passent la correction de Holm dans la comparaison principale
(\cref{tab:sondage-pineau,tab:sondage-extension}). Le code et le calcul résistent à toutes les coupes
que nous avons tentées, la sensibilité à presque toutes. Le code garde son verdict quand les deux
années sont lues sur la même sorte de version, actes ou arXiv (\cref{sec:sr-robustesse}).
L'incertitude rapportée est le changement le plus fragile : elle ne passe plus la correction quand
on retire les deux articles de la comparaison lus en partie par un autre modèle. Son verdict dépend
donc de deux articles. Les lecteurs de 2018 n'étant pas ceux de 2026, huit lecteurs neufs ont relu
les deux années : les hausses du code, du calcul et de la sensibilité tiennent, celle de
l'incertitude rapportée garde son sens sans être établie (\cref{sec:sr-annees}). Un lecteur d'une
autre famille de modèles retrouve le sens des quatre hausses. Quand ses \enquote{oui} sans citation
retrouvée comptent \enquote{non}, il ne confirme pas celle de l'incertitude rapportée
(\cref{sec:sr-gemma}). Le critère X8 montre ce que mesure cette dernière hausse.
La part des articles qui calculent leur incertitude sur des graines aléatoires passe de 35~\% à
42~\%, un écart que le sondage ne détecte pas : la barre d'erreur s'est répandue plus vite que la
barre calculée sur les entraînements indépendants (\cref{sec:sr-principale}). Lues sur 20 articles
par année entre les deux dates, les quatre pratiques montent lentement, sans saut juste après la
liste de contrôle que NeurIPS impose depuis 2019 (\cref{tab:trajectoire}).

\paragraph{Le jugement du résultat a peu bougé} Le réglage des références stagne : hyperparamètres des références donnés (Q4), même soin de réglage
annoncé (X4) et origine de leur code (Y5). La version des bancs d'essai (Y7), la définition de la
mesure de performance (Y1) et la sélection des meilleurs essais (X3) restent constantes. Le test
statistique, l'évaluation tenue à part, les hyperparamètres de la méthode et la façon de les choisir
montent de 6 à 22 points selon le critère, aucune de ces hausses ne passant la correction. Avec une
puissance de 80~\%, le sondage ne pouvait détecter qu'une hausse de 18 à 30 points selon le critère,
et de 26 à 39 après correction (\cref{sec:val-stat}) : il ne dit pas que ces pratiques stagnent, il
ne peut pas le savoir.

\paragraph{Une case cochée ne garantit pas la pratique} Les 18 articles de NeurIPS 2025 qui répondent oui à la
question de leur liste de contrôle sur la signification statistique n'appliquent aucun test. Sur le
code et sur le calcul, une partie des \enquote{oui} n'a ni lien ni chiffre derrière elle
(\cref{sec:sr-declarations}).

\paragraph{Le hasard trompe l'œil, le tri trompe le test} Deux groupes de cinq essais tirés au hasard
parmi ceux d'un même programme trompent rarement un test statistique. Triés, les cinq meilleurs
contre les cinq pires le trompent de 26~\% à près de 100~\% des fois selon le couple (programme,
tâche), 31~\% dans le cas principal, PPO sur \task{HalfCheetah} (\cref{sec:e1}). Garder les trois
meilleurs de dix essais gonfle la moyenne de 32~\% en médiane sur les seize couples
(\cref{tab:e4}). Le $t$ de 9,09 de 2018 s'explique probablement par le script publié par Peter
Henderson, qui mesure la dispersion des épisodes et ignore celle des graines aléatoires. La paire
($t$, $p$) publiée exclut les autres lectures du calcul plus qu'elle ne désigne ce script : la
lecture concurrente, la moyenne des $p$, n'a pas de degrés de liberté bien définis. Appliqué à nos
essais de TRPO, ce script dépasse 9 dans 38~\% des tirages au hasard (\cref{tab:b1}).

\paragraph{Le meilleur programme dépend de la tâche, du budget et de la physique} Sur nos quatre programmes, le meilleur change avec la tâche.
PPO y dispose pourtant de cinq fois plus de pas que SAC et CrossQ. À budget égal et avec vingt
graines aléatoires, CrossQ mène sur les quatre tâches, mais l'ordre de deux paires s'inverse encore
de façon établie d'une tâche à l'autre. L'inversion de PPO JAX et de SAC résiste même à la correction de Bonferroni (\cref{sec:e2}).
Les deux écritures de PPO ne diffèrent de façon établie que sur \task{Swimmer}, d'un écart petit en
points de score (\cref{sec:e3}). Sous une physique modifiée de 20~\%, CrossQ bat chacun des trois
autres programmes sur \task{Hopper}, à dix comme à vingt graines aléatoires. PPO JAX y perd contre SAC, qu'il battait en physique d'entraînement (\cref{sec:e7}).

\paragraph{Une recette publiée a vieilli} Les agents entraînés avec la recette publiée par RL Baselines3 Zoo ne font que se tenir debout sur
\task{Hopper}. L'agent que le Zoo a lui-même publié en 2022, avec les mêmes réglages mais une pile
logicielle plus ancienne, avance et atteint un score voisin de celui de nos agents par défaut. La
recette fonctionnait avec la pile d'origine et a vieilli (\cref{sec:e6}).

\paragraph{Sur les modèles de langue, aucun gain de GRPO n'est établi} Avec cinq graines aléatoires, aucun gain à l'intérieur
d'une configuration ne peut être établi, puisque le test exact de signe ne descend pas sous 0,0625.
Un contraste est établi : chez Llama, la vraie récompense fait mieux que la récompense tirée au sort,
parce que celle-ci détruit le modèle. Sous la vraie récompense, Llama perd en justesse gloutonne et en
pass@8 dans chacune des cinq graines aléatoires, alors que son pass@1 monte un peu. Chez Qwen, le gain
de pass@1 s'explique par la fin des réponses coupées : la justesse des réponses finies ne monte pas.
Qwen a de plus tourné avec une pénalité de répétition que Llama n'avait pas
(\cref{chap:07-modeles-de-langue}). 

\section{Recommandations}
\label{sec:recommandations}

Les recommandations suivantes découlent des mesures de ce rapport. Chacune nomme la faille qu'elle
vise et le résultat qui la justifie.

\subsection{Pour les auteurs}

\begin{enumerate}
  \item \textbf{Annoncer le nombre de graines aléatoires et en lancer au moins dix.} Sur 41 articles
        de RL expérimentaux du premier tirage de 2026, 20 ne disent pas combien d'essais ils
        lancent. La médiane des autres vaut cinq. Avec cinq graines aléatoires, un écart de moins de
        20~\% se lit à l'envers près d'une fois sur trois, plus d'une fois sur trois quand vingt
        graines aléatoires fixent le bon ordre ou quand l'ordre se juge hors échantillon (E5,
        failles F9 et F10).
  \item \textbf{Ne jamais trier ni sélectionner les essais.} Rapporter tous les essais, y compris
        ceux qui échouent. Dire ceux qui ont été exclus (E1, E4, critères X3 et Y6).
  \item \textbf{Tester sur les graines aléatoires.} L'entraînement complet est l'unité indépendante. Un
        test qui dénombre les épisodes déclare une différence entre deux groupes du même code (B1, faille
        F11). C'est la même chose pour la barre d'erreur : moins d'un article expérimental sur deux paru
        en 2026 la calcule à partir de graines aléatoires (X8).
  \item \textbf{Montrer un intervalle à côté de chaque courbe et dire sur quoi il est calculé.}
        Nommer l'écart type, l'erreur type ou l'intervalle de confiance, puis la dispersion qu'il
        résume : 14 des 44 articles de 2026 qui tracent une barre ne disent pas sur quoi elle est
        calculée (critères X2, X8, Y2, Y11).
  \item \textbf{Comparer à budget égal, ou dire l'écart de budget.} Un classement à budgets
        inégaux compare aussi les budgets (E2 à budget égal).
  \item \textbf{Corriger pour les comparaisons multiples et dire ce que l'échantillon pouvait
        voir.} Quatre tâches testées sans correction ont fait lire à E3 deux différences, dont une
        seule résiste à la correction de Holm. Un \enquote{pas de différence} se lit avec la plus
        petite différence que l'échantillon détecte (E3, \cref{sec:val-stat}).
  \item \textbf{Régler les références avec le même soin que la méthode et le dire.} Donner le
        budget de réglage de chacune, ses hyperparamètres finaux et la version des bibliothèques où
        sa recette a été réglée (E6, critères Q4, X4, Y5).
  \item \textbf{Évaluer dans des conditions tenues à part.} Graines aléatoires d'environnement
        jamais vues, dynamiques modifiées ou distribution de test distincte (E7, critères Q7, Y8).
  \item \textbf{Donner les versions.} Celle de chaque banc d'essai et de chaque dépendance, dans un
        fichier d'environnement figé (critères Y7, Y10).
  \item \textbf{Pour un modèle de langue : plusieurs familles de modèles, tout le réglage du
        décodage et plusieurs mesures.} Un résultat obtenu sur une seule famille renseigne sur cette
        famille. Un réglage hérité du fichier de configuration du modèle, comme une pénalité de
        répétition, change ce que l'on échantillonne. Un gain de justesse se lit avec la limite de
        longueur, la part coupée avant et après, la justesse des réponses complètes, la réponse
        gloutonne et pass@$k$ (\cref{chap:07-modeles-de-langue}, critère X6).
  \item \textbf{Écrire le protocole avant le premier essai et publier chaque changement.} Dire,
        pour chaque changement, ce qui avait déjà été vu au moment de le décider.
  \item \textbf{Recalculer chaque nombre publié depuis ses fichiers.} Un script qui confronte chaque
        nombre du texte à son fichier de résultats coûte peu et attrape les dérives
        (\cref{sec:verification}).
\end{enumerate}

\subsection{Pour les relecteurs}

Les sept questions de l'article de blog s'adressent à tout lecteur. Un relecteur peut les poser à
chaque soumission.
\begin{enumerate}
  \item Combien de graines aléatoires ? Avec moins de cinq essais, ou avec \enquote{les meilleurs
        essais}, le résultat dit peu de chose.
  \item Quelle incertitude ? Une moyenne sans intervalle ni bande autour de la courbe ne permet pas
        de juger un écart.
  \item Quelle méthode de référence, réglée comment ? Si la méthode a reçu un mois de réglages et la
        référence ses valeurs par défaut, l'écart mesure le mois de réglages.
  \item Quel budget ? Comparer deux méthodes entraînées sur des nombres de pas différents, c'est
        aussi comparer ces nombres de pas.
  \item Évalué où ? Un agent jugé sur les situations de son entraînement a pu les apprendre par cœur.
  \item Quel code ? Sans code publié, le lecteur ignore ce que l'article a vraiment fait.
  \item Quel calcul ? Un résultat que personne d'autre n'a les moyens de recalculer ne se reproduit
        pas.
\end{enumerate}
Le relecteur peut ajouter une huitième question, que ce rapport justifie : sur quelle unité
l'incertitude et le test sont-ils calculés ? Une barre ou un test sur les épisodes, ou sur les points
d'une courbe empilés, ne mesure pas la variabilité entre entraînements.

\subsection{Pour les conférences}

\begin{enumerate}
  \item \textbf{Demander le test en plus des barres.} La question de NeurIPS sur la
        signification statistique accepte des barres d'erreur \enquote{ou une autre information
        appropriée}. Les 18 articles de notre tirage qui y répondent oui n'appliquent aucun test. Une question qui
        demande l'unité du test et le nombre d'entraînements indépendants serait plus exigeante.
  \item \textbf{Vérifier une case cochée sur un échantillon.} Certains articles répondent oui sur le code
        sans lien, ou sur le calcul sans chiffre. Un tirage au sort d'une fraction des articles acceptés
        mesurerait l'écart entre ce qui est déclaré et ce qui est pratiqué.
  \item \textbf{Exiger le lien vers le code dans la version finale.} \enquote{Le code sera publié}
        ne vaut pas un lien. Un lien ajouté à la seule version arXiv échappe aux actes.
  \item \textbf{Demander la version de chaque banc d'essai.} Elle n'a pas bougé de façon détectable depuis
        2018 (23~\% puis 26~\%), quand les trois modèles d'IA agentique qui notent le coût de chaque
        pratique (\cref{chap:04-sondage-resultats}) estiment qu'elle est presque gratuite.
  \item \textbf{Soutenir les défis de reproduction.} Le défi d'ICLR 2018 présenté par Pineau a montré
        qu'une reproduction par des tiers est faisable : parmi les participants qui ont répondu au
        sondage, 54,5~\% en ont reproduit une partie, 32,7~\% l'essentiel
        \parencite[diapositives 37 à 43]{pineau2018iclr}. Le refaire par grandes familles de
        résultats, avec un protocole écrit d'avance, mesurerait ce que la liste de contrôle ne
        mesure pas.
\end{enumerate}

\section{Travaux futurs}
\label{sec:futurs}

\begin{itemize}
  \item \textbf{ICLR dans la trajectoire.} Les années 2019 à 2024 manquent d'ICLR. Une lecture de
        ses actes, par une voie de téléchargement qu'OpenReview accepte, compléterait la série.
  \item \textbf{Un sondage plus grand sur les critères qui montent sans être établis.} Le test
        statistique, l'évaluation tenue à part et la méthode de choix des réglages montent de 6 à 22
        points sans que le sondage puisse l'établir. Un échantillon assez grand pour détecter une
        hausse de dix points trancherait.
  \item \textbf{Une lecture, avec la même grille, des articles qui appliquent ce qu'ils appellent RL
        à de grands modèles de langue}, sur un échantillon de taille comparable, dirait si
        l'évaluation tenue à part, qui y progresse, s'accompagne des autres pratiques.
  \item \textbf{Une relecture humaine à l'aveugle} de trente articles, quinze par année, mesurerait l'écart
        entre les agents et un lecteur humain.
  \item \textbf{Un nouvel entraînement complet sur une autre plateforme}, Linux sur x86-64,
        testerait les statistiques de nos expériences avec d'autres tirages (\cref{sec:reentrainer}).
  \item \textbf{D'autres familles de tâches} que la locomotion MuJoCo : jeux à images, RL hors
        ligne, robotique réelle. Il reste à vérifier si les conclusions de E2 à E7 s'y
        généralisent.
  \item \textbf{Le volet GRPO à plus grande échelle.} Si une relance à dix graines aléatoires établissait son contraste, le même protocole sur des modèles de 7 milliards de paramètres rapprocherait le
        volet de \textcite{shao2025spurious}.
\end{itemize}

\paragraph{Ce que nous avons corrigé}
Pineau demandait en 2018 d'être aussi critique envers ses bons résultats qu'envers ses mauvais.
Nous avons cherché nos propres erreurs et corrigé chacune. Dans E6, notre code donnait à la
normalisation des récompenses du Zoo un facteur d'actualisation de 0,99 au lieu de 0,999 : la recette
corrigée ne fait pas mieux de façon détectable (\cref{sec:e6}). Le décompte de E2 comptait des triplets
au lieu des paires qui s'inversent : le décompte strict trouve trois paires au lieu de quatorze
changements et le verdict tient (\cref{sec:e2}). Le facteur qui relie le $t$ du script de Henderson à
celui d'un test sur les graines aléatoires vaut $\sqrt2/\rho$ pour nos dix épisodes, non $1/\rho$
(\cref{sec:b1}). Dans le volet GRPO, le gain de Qwen, net à trois graines aléatoires, ne l'est plus à
cinq : il tient à la limite de longueur, qui coupait ses réponses. Qwen décodait aussi avec une
pénalité de répétition que Llama n'avait pas et le pilote qui a fixé la tâche avait lu le jeu de test :
les chiffres publiés restent les mêmes, la comparaison entre familles mêle la famille et le décodage
(\cref{chap:07-modeles-de-langue}). Le dépôt garde la trace de chaque correction pour qui voudra
refaire le travail.

\section{Ouverture}\label{sec:ouverture}
Michael I. Jordan décrivait dès 2019 l'intelligence artificielle comme une discipline d'ingénieurs
encore à bâtir, plus que comme une révolution déjà faite
\parencite{mijordan2019revolution,mijordan2019drai}. Il y revient à propos d'un incident récent où des
agents IA ont été laissés libres d'agir : ce sont, dit-il, \enquote{surtout des mauvaises pratiques
d'ingénierie}. Il compare : \enquote{en ingénierie chimique, on n'aurait jamais eu une situation où
l'on aurait laissé des explosions se produire pour voir jusqu'où les choses pouvaient mal tourner}
\parencite{mijordan2026liberation}.

Les trois R de Pineau disent la même exigence à l'échelle d'un article. Un résultat reproductible,
des pièces réutilisables, une conclusion robuste : c'est ce qu'un ingénieur attend de tout ouvrage
avant de s'en servir. Notre sondage montre que le RL a adopté les marques extérieures de cette
discipline, le code publié, le calcul déclaré, la barre d'erreur. La façon de juger un résultat a
beaucoup moins changé : le nombre d'essais, le test statistique, l'évaluation tenue à part.

Michael I. Jordan voit aussi l'IA comme un phénomène économique et social. Il y a des producteurs,
des consommateurs et ceux qui décident dans quels domaines ils l'acceptent. Pour les programmes IA
qui apprennent par essais et erreurs, leur confiance passera par des résultats que d'autres peuvent
refaire. Les méthodes existent. Un protocole écrit d'avance coûte peu à côté d'une conclusion fausse.
